\section{Q-KM：长期价值、拒单与真实调度代价}
\subsection{用区域共享价值连接当前与未来}
贪心最近优先短接驾，KM即时收入优先当前高车费；后者可能选择把司机带到远离后续需求的长途订单，司机忙时和新位置都会影响全天收入。历史学习区域价值、再做组合派单的思路参照Xu等\cite{xu}，本次采用教师的表格化教学近似，并非复现其生产平台。状态$(t,i)$包含时间片和8×8区域，动作是去目的地区域$k$或WAIT，108×64×65的Q表由所有司机共享。它节省了逐车状态空间，但损失同区域内具体位置、当前订单竞争和司机供给分布的信息，不能假定同区域的车具有相同接单机会。

一条转移持续$\Delta$片，其TD目标为
\begin{equation}
y=r+\gamma^{\Delta}\max_{a'\in A(t',i')\cup\{WAIT\}}Q(t',i',a'),
\end{equation}
到达一天结束之后时bootstrap项为0。$A$必须是该司机下次可用时3 km内实际活动订单的目的地区域，不能对所有65动作盲目取max。访问计数给出$\alpha=\max(1/N,0.05)$，同时以可行动作max样本更新区域$V$；经验按时间逆序回放。教师原始手算断言与终止/WAIT反例均实际执行。

Q-KM把边权写为
\begin{align}
w_{dj}&=fare_j+\beta LT(t,i_d,k_j),\\
LT(t,i,k)&=\gamma^{\Delta(i,k)}V(\min(t+\Delta,T),k)-\gamma V(t+1,i).
\end{align}
区域中心近似$\Delta(i,k)$用于价值查表，实际忙时仍用该车的接驾加送客距离、25 km/h与10 min片长向上取整。不可行和非正边不派，$\beta=0$恰为KM即时收入。忙车在释放前不能再接单；位置在释放时更新。订单允许年龄0/1/2片，超过2片取消；已接、期末仍未分配、取消和期末服务未完成分别计数。

收入保持教师语义，在接受订单时确认车费，可能包含horizon后才完成的服务；另报告horizon内已完成收入，避免把接单与完成混为一谈。训练每回合1500个重采样代理订单、共300实际模拟回合，$\gamma=.92$和$\beta=.5$是教学参考而非已知最优值。

\datafigure{rl_learning}{60车、开发训练seed2024的实际300回合。左为3个固定验证模拟日的接单车费代理，基线同日同初始条件；右为逐回合TD绝对误差，20回合滚动均值仅辅助阅读。}{fig:learning}
\datafigure{rl_value_maps}{训练来源的区域价值地图，所有时段共用色标。橙色圆面积为同一批3000训练代理起点的8×8网格计数，不是人口/真实出租车需求；价值是估计量。}{fig:vmap}

\subsection{beta、车队规模与控制变量}
验证扫描$\beta=0,.2,.4,.5,.6,.8,1$，同时记录被非正边权拒绝的可行边比例。60车开发模型在$\beta=1$时收入下降、拒单增多；这提供过估计或机会成本失配的机制线索，但不能仅凭收入曲线证明Q值偏差。GAMMA=.8/.95、GRID=6/10、PATIENCE=1/3分别独立重训300回合。增加PATIENCE改变乘客等待合同，因此即使收入较高也不能无条件当同服务标准下的改进。

\datafigure{rl_beta_rejection}{同一训练模型、5个验证模拟日，beta变化的收入与非正可行边比例。未使用最终测试曲线选参数；边比例不是订单取消率。}{fig:beta}
\begin{table}[htbp]\centering\small
\caption{60车开发单因素；每项重新训练300回合、同5个验证模拟日。}\label{tab:rlsensitivity}
\begin{tabularx}{\linewidth}{Yrrr}\toprule
变化 & 接单车费 & 应答率 & 接驾/km\\\midrule
原参考 & 16,830.68 & 0.5900 & 1.891\\
gamma=.8 & 16,089.53 & 0.5345 & 1.894\\
gamma=.95 & 16,574.42 & 0.6085 & 1.913\\
grid=6 & 16,748.56 & 0.5784 & 1.888\\
grid=10 & 16,761.73 & 0.5937 & 1.910\\
patience=1 & 16,366.49 & 0.5747 & 1.880\\
patience=3 & 17,290.17 & 0.6064 & 1.890\\
\bottomrule\end{tabularx}
\tabnote{每项只改变命名参数；patience变化改变等待合同，不能直接解释为相同服务约束下的质量改善。}
\end{table}


最终确认在前70\%日期完整重训，对最后30\%日期池的5个固定模拟日评价。表\ref{tab:rlfleet}给出三车队规模，同一车队的策略共用订单和初始司机。40车的长期价值帮助更明显，60车Q-KM相对KM收入增加2.47\%，100车反而下降3.59\%。司机充足时继续拒绝即时可行单的代价可能超过未来价值收益；这里的有限模拟证据不能变成普适车队阈值。

\begin{table}[htbp]\centering\small
\caption{首批5个留出模拟日：车队和策略比较。}\label{tab:rlfleet}
\begin{tabularx}{\linewidth}{rlrrrr}\toprule
车辆 & 策略 & 接单车费 & 应答率 & 接驾/km & 期内完成车费\\\midrule
40 & 最近 & 13,333.15 & 0.5205 & 1.660 & 13,308.94\\
40 & KM即时 & 12,437.18 & 0.4069 & 1.901 & 12,428.08\\
40 & Q-KM & 13,994.90 & 0.5169 & 1.916 & 13,985.80\\
60 & 最近 & 17,102.75 & 0.6749 & 1.530 & 17,078.54\\
60 & KM即时 & 16,757.42 & 0.6072 & 1.907 & 16,744.18\\
60 & Q-KM & 17,170.85 & 0.6603 & 1.928 & 17,161.74\\
100 & 最近 & 20,716.17 & 0.8155 & 1.291 & 20,712.03\\
100 & KM即时 & 20,793.01 & 0.8080 & 1.912 & 20,773.72\\
100 & Q-KM & 20,046.49 & 0.7795 & 1.953 & 20,037.38\\
\bottomrule\end{tabularx}
\tabnote{接单收入与期内完成收入分开。车费为教学代理，非真实人民币；各车队独立重训300回合。}
\end{table}

\datafigure{rl_fleet_final}{最终后段日期池、每车队5个相同采样模拟日。各车队独立重训300回合；均值和波动按模拟日计算，未把它们称为5天真实运营。}{fig:fleet}

\subsection{Double Q是否修复beta=1}
Double Q以一张表在真实可行动作上选argmax，另一张表评价该动作，随机更新其中一表；分别记录访问次数。独立手算核验了动作选择、交叉评价、持续时间折扣与终止边界。初次实现的更新coin与行为探索共用随机流，会改变同seed下后续探索；停止受影响运行、保存旧结果，再分离更新RNG。Q-only的Q/V/计数/曲线/TD/经验数组逐项bit等价，旧Double不用于正式公平对照。

五训练seed、每seed五个验证模拟日中，Double在$\beta=1$提高收入，却降低应答，非正边比例也未降低。区域V与已实现策略折扣收益之间的描述性gap减小，但二者访问分布和策略不同，不能把该gap当成已知最优值的无偏估计误差。Double降低统计max偏差的思想\cite{double}也不能直接消除多个司机竞争同一订单的结构性失配。

\datafigure{rl_double_paired}{Q与Double在beta1的成对训练seed比较。收入、应答和非正可行边比例分别展示；每点为对应训练seed下固定验证模拟日均值，V诊断另由实际访问样本记录。}{fig:double}
在最终5日确认中，Double的收入14,336.25高于Q的13,995.18，应答率却由0.6029降至0.5728，而且两者均低于$\beta=.5$参考。因此完成了加分的真实机制检验，但不把它升级为全面优胜方案。

\subsection{主动调度：先付出移动成本，再谈收益}
对没有匹配的空闲司机，先比较等待价值与四邻域中心的折扣到达价值。移动占用真实向上取整的行驶片数，期间不可接单，释放后才在新位置派单；额外空驶距离和时间单独累积。原始无成本门槛会产生大量调度，60车开发模拟日额外空驶约5000 km，不能只报告gross收入增加。

由这一负结果触发单因素移动门槛：在候选移动价值中扣$\lambda\,d_{move}$，比较$\lambda=0,1,3,5$，其单位为教学价值/公里。它是控制移动强度的教学参数，不是已知真实油电价格。先保持原Q-KM边权逐项比较，再比较移除LT边权的KM父项。五训练seed验证中，KM即时边权配合同一个learned V调度在全部五seed的gross收入高于Q-KM+同调度，说明两项机制未必互补；最终冻结回退为KM即时边权+$\lambda=3$调度，保留Q-KM父项和$\lambda=5$低空驶替代。

\datafigure{rl_reposition_cost}{同一验证输入的移动门槛与收入/额外空驶权衡。extra empty仅指reposition，接驾距离另记；成本线是预登记情景，不能当真实利润。}{fig:move}
\datafigure{rl_movement_ablation}{训练seed2024、5个成对验证模拟日的组件与移除比较：相同learned V用于移动，派单边权是否仍含LT为单独因素。五训练seed确认另见正文取舍，避免把该图误读为五次独立训练。}{fig:ablation}
\begin{table}[htbp]\centering\small
\caption{冻结策略的首批5日留出确认，60辆车。}\label{tab:rlenhance}
\begin{tabularx}{\linewidth}{Yrrrr}\toprule
策略 & 接单车费 & 应答率 & 接驾/km & 额外调度/km\\\midrule
Q-KM参考 & 17,170.85 & 0.6603 & 1.928 & 0.00\\
KM父项 & 16,757.42 & 0.6072 & 1.907 & 0.00\\
KM+调度3 & 17,811.65 & 0.6861 & 1.929 & 826.55\\
Q-KM+调度3 & 17,280.52 & 0.6817 & 1.931 & 685.44\\
Q-KM+调度5 & 17,573.25 & 0.6875 & 1.942 & 278.01\\
Q(beta1) & 13,995.18 & 0.6029 & 1.960 & 0.00\\
Double(beta1) & 14,336.25 & 0.5728 & 1.945 & 0.00\\
融合膝点 & 14,750.38 & 0.5396 & 1.053 & 0.00\\
融合约束 & 16,932.91 & 0.6493 & 1.360 & 0.00\\
\bottomrule\end{tabularx}
\tabnote{额外调度距离不包含接驾；所有策略同订单/初始车。原Q、Double、融合和父项分别保留，不据此更换选择。}
\end{table}


首批5个留出模拟日，冻结KM+调度相对Q-KM的gross收入差640.81，描述性t区间为$[-64.07,1345.68]$，跨0；每模拟日额外调度826.55 km。随后只对已冻结策略增加20个新采样seed，配对差435.45、区间$[141.68,729.23]$，额外848.08 km。两批分别保留，第二批仅减少同一代理池的Monte Carlo不确定性，既不是20个新真实日期，也没有用于重选参数。

即便只假设每公里额外调度成本为1，补充20seed的收入减额外调度成本之差为$-412.62$，描述性t区间$[-717.03,-108.22]$；这一诊断还没有扣接驾、载客和固定成本，故不能称真实净利润。更省空驶的替代方案与gross最佳方案是权衡，不据最终读数重新命名胜者。
\begin{table}[htbp]\centering\small
\caption{固定KM+调度3相对Q-KM：两批模拟身份保持独立。}\label{tab:rlmc}
\begin{tabularx}{\linewidth}{lrYrr}\toprule
批次 & n & 比较量 & 配对均差 & 描述性t95区间\\\midrule
首批 & 5 & gross收入差 & 640.81 & [-64.07, 1,345.68]\\
首批 & 5 & 收入减额外调度成本差 & -185.74 & [-888.44, 516.95]\\
补充 & 20 & gross收入差 & 435.45 & [141.68, 729.23]\\
补充 & 20 & 收入减额外调度成本差 & -412.62 & [-717.03, -108.22]\\
\bottomrule\end{tabularx}
\tabnote{成本情景只扣1教学单位/km的额外reposition，没有扣接驾/载客/固定成本；不是实际利润。区间仅描述相同代理池的采样波动。}
\end{table}

\datafigure{rl_final_tradeoffs}{首批5个留出模拟日的成对收入差与额外调度成本。补充20个新采样模拟日另见表；区间描述代理池内的采样波动，不覆盖真实业务不确定性。}{fig:finaltrade}

\subsection{Q-KM加NSGA-II：估计前沿不等于全天收入}
完成Q-KM和独立NSGA后，把每条可行正边的$fare+\beta LT$与接驾km作为两个目标，用全部可行候选边进行NSGA-II。动态每次决策采用种群40、30代，种子和选择规则在验证时固定；这与Part A的100/300科研比较分开记录实际预算。膝点规则和接驾不超过同状态Q-KM总接驾60\%的约束规则分别执行，估计边权目标、前沿和实际全日模拟均保存。

最终约束方案把平均接驾从Q-KM的1.928 km降至1.360 km，但收入由17,170.85降至16,932.91，应答也下降。膝点取更低接驾时收入/应答代价更大。它实现了真实的双目标权衡，却没有证明比单独Q-KM全面更好，所以保留为教师加分成果而不强制采用。
\datafigure{rl_fusion_front_and_effect}{验证阶段：左侧为实际状态的静态100种群/300代估计前沿，右侧为动态40种群/30代策略的完整验证模拟效果；上文留出数值另表报告。不用估计边权替代已实现收入。}{fig:fusion}
